% conclusion.tex
%
% Aetf <aetf@unlimitedcodeworks.xyz>
% Copyright 2016 Aetf <aetf@unlimitedcodeworks.xyz>
%
% multiple1902 <multiple1902@gmail.com>
% Copyright 2011~2012, multiple1902 (Weisi Dai)
%
% Project Home: https://github.com/Aetf/xjtuthesis
%
% It is strongly recommended that you read documentations located at
%   https://github.com/Aetf/xjtuthesis/wiki
% in advance of your compilation if you have not read them before.
%
% This work may be distributed and/or modified under the
% conditions of the LaTeX Project Public License, either version 1.3
% of this license or (at your option) any later version.
% The latest version of this license is in
%   http://www.latex-project.org/lppl.txt
% and version 1.3 or later is part of all distributions of LaTeX
% version 2005/12/01 or later.
%
% This work has the LPPL maintenance status `maintained'.
%
% The Current Maintainer of this work is Aetf.
%
\chapter{结论与展望}
\echapter{Conclusions and Future Work}

\section{全文工作总结}
\esection{Summary of This Thesis}

本文围绕基于信息交互的多智能体强化学习算法展开研究，重点关注 CTDE 框架下信息交互机制的设计、失效模式及改进方向。针对现有通信方法在消息语义、通信效率和融合策略方面存在的问题，本文以 MAPPO 为主干算法，首先构建了相对可靠的无通信主干基线，然后在其基础上系统探索了多种轻量通信结构，并对攻击通信流进行了系统性的机制修复。

在方法层面，本文形成了一个较为清晰的设计思路：通信模块不应简单扩大消息容量或替代整个局部策略，而应以轻量、残差、受约束和面向协作子空间的方式接入。围绕这一思路，本文先后实现并比较了最小侵入通信、锐化路由通信、动作意图共享、攻击子空间融合以及攻移双流通信等结构。在此基础上，针对攻击通信流注意力退化这一关键问题，本文进一步引入选择性熵正则、沉默预算和意图对齐等机制修正，使跨种子平均峰值胜率提升至 $0.746$。随后，本文又将研究重点从“通信是否存在”推进到“通信能否稳定进入动作层，并在更需要协同的状态中介入”，最终形成了稳定通信介入方案。进一步的边界诊断表明，继续平滑对齐样本选择、单纯提高通信门控强度，或全局性鼓励通信改变动作 logits，均不能稳定增加执行期通信贡献，后续改进需要转向状态条件化的通信因果性建模。

在实验层面，\texttt{join1} 环境验证了通信机制并非天然有效，设计不当时可能导致完全坍塌。在 SMAC \texttt{5m\_vs\_6m} 中，本文系统分析了主干训练的崩溃-恢复动力学和种子敏感性，揭示了冷启动通信的结构性失败原因，包括门控与学习率的同构性、通信预热的不可持续性，以及隐藏状态融合仅改变表示却未能稳定转化为胜率提升的局限。通信隔离评估进一步表明，即便机制基本正确，攻击通信也并非对所有种子都稳定有益：它在部分种子上产生了约 $+5\%$ 的真实边际贡献，但在其他种子上无显著影响。进一步的同伴目标与本地目标因果诊断显示，早期阶段的问题不是同伴信号不存在，而是有效动作扰动近乎为零；在后续修正中，通信已经能够真实进入攻击动作层，并最终在稳定通信介入方案的代表性检查点上达到更稳、更健康的工作区间。该检查点同时具备较强性能、健康的双流通信参数和清晰的机制边界：攻击流几乎不再通过“无通信”标记逃逸，而移动流则保持较高的按需沉默比例，从而共同形成较低的总体执行期通信预算。进一步的门控下限、Huber 杠杆与更激进的因果修正实验还表明，通信变强、甚至通信确实改变动作 logits，都不等价于通信变好；若缺少因果状态选择，增大通信注入强度反而可能破坏已有策略。

除主实验地图外，本文还在 \texttt{MMM2} 与 \texttt{2c\_vs\_64zg} 上完成了跨地图补充验证。两组结果共同说明，当前轻量双流通信框架并非仅在 \texttt{5m\_vs\_6m} 上偶然有效，而是在不同复杂协作结构下都能够表现出正向迁移收益。其中，\texttt{MMM2} 的结果表明通信有效性与主干能力阈值高度相关：当通信在较弱主干阶段过早介入时，并不优于同预算的纯主干续跑；而当主干已进入更强协作区间后，再接入通信则能够反超同预算主干基线。\texttt{2c\_vs\_64zg} 则给出了更干净的正例：从相同强主干检查点出发，通信续跑最终达到 $0.7625$ 的测试胜率和 $20.1492$ 的平均回报，显著优于纯主干续跑的 $0.5625$ 胜率和 $19.4775$ 回报。这进一步强化了本文的总体叙事：通信收益是可迁移的，但其释放方式依赖于任务结构与主干所处的能力阶段。

\section{主要结论}
\esection{Main Findings}

结合全文研究，本文的主要结论可以概括为以下几点。

第一，构建强基线比盲目堆叠通信结构更重要。主干策略本身的训练动力学极其复杂，存在崩溃-恢复循环和种子间的分叉行为。若主干本身不稳定，则通信实验的结论将无法归因。

第二，通信在冷启动下结构性失败。IGG、固定门控、通信预热和隐藏状态融合等策略均无法使通信在从零训练中产生稳定正向收益。门控值的功能等价于通信学习率，这一洞察直接催生了门控退火策略和后续的两阶段训练设计。

第三，通信内容的语义清晰度比通信公式的复杂程度更重要。动作意图之所以有效，关键在于其直接对应协作任务中的集火需求。当注意力机制将消息路由到正确的队友后，通信才开始产生真实价值。

第四，通信路由的稳定化是机制修复的核心。攻击注意力在引入选择性熵正则和沉默预算后，最终能够在 entropy $0.40$--$0.45$ 的理想区间内稳定运行。通信隔离评估首次提供了攻击通信有效性的因果证据：在 seed 3 上可测量到 $+4.7\%$ 的执行期正向增益。

第五，主干质量是通信有效性的前提条件。实验中发现，若主干策略在通信热启动微调前未达到足够稳定的状态，通信支路会放大优化噪声并导致主干退化。这一发现将通信研究的焦点从“设计更好的通信”转向“确保主干质量足以支撑通信训练”，这也是两阶段训练策略的核心动机。

第六，通信触达动作并不等价于通信具有因果收益。早期的同伴目标与本地目标因果诊断已经显示：同伴意图存在、冲突状态非零，但最终融合策略几乎不跟随同伴目标；后续修正则进一步表明，通信可以被稳定推进到攻击动作层，并以可测量的有效动作扰动影响动作 logits。特别是稳定通信介入方案的代表性检查点展示了当前最干净的方案状态：攻击流的无通信概率近零，移动流保持较高按需沉默，攻击门控处于中等区间，有效动作扰动稳定为正，注意力熵保持健康，同时 final win 达到 $0.9375$。从总体执行期预算看，该检查点的双流通信成本约为 $21.36$，仅约为 repo MAIC 的 $44.5\%$。然而，即使在这一检查点中，融合策略保持本地目标的比例仍接近 $0.98$，说明通信虽然已经在动作层面稳定工作，却仍未频繁跨过本地决策间隔。

第七，更强的动作层暴露本身不是答案。门控下限实验将攻击门控从约 $0.071$ 提升至 $0.122$，有效 logit 扰动也约翻倍，但三 seed 最终平均胜率从 $0.667$ 下降至 $0.479$。进一步的 Huber 优势引导杠杆虽将攻击门控提升至 $0.171$，并将有效扰动提升至约 $4.95\times10^{-3}$，但 128 局通信隔离评估中 normal $-$ no\_attack 的平均差值仍为 $-0.52\%$。这说明当前瓶颈不是通信没有被打开，也不只是通信没有改变动作，而是通信尚未学会在真正需要干预的状态中沿有利方向修正本地策略。换言之，本文已经较好地修复了通信进入动作层的暴露瓶颈，但因果性决策改变瓶颈仍未彻底解决。

第八，跨地图结果支持“主图成立、复杂图可迁移”的总体叙事，但不支持无条件强泛化。本文在 \texttt{MMM2} 与 \texttt{2c\_vs\_64zg} 上都观察到通信相对于强主干策略的正向收益，其中前者强调能力阈值敏感性，后者强调在强主干条件下的稳定额外增益。因此，当前更准确的结论不是“通信在所有地图上都稳定强于主干”，而是“在主图上方法成立，在复杂补充地图上也观察到正向迁移，且通信有效性与任务结构、主干阶段和介入时机密切相关”。

\section{后续工作展望}
\esection{Future Work}

尽管本文已取得若干阶段性结论，但仍有多个方向值得继续深入。

首先，本文当前的 hidden-state fusion 验证已证明通信能够改变策略表示（cosine 下降至 $0.816$），但并未将其转化为胜率提升。后续可尝试在 policy head 中引入非线性处理（如 MLP head 替代当前的单层线性 head），以充分利用通信提供的正交方向信息。

其次，本文发现主干与通信联合训练的动力学稳定性是当前最核心的瓶颈。后续可尝试引入主干保持正则，例如在通信热启动微调阶段对局部策略施加 KL 约束，确保主干在通信训练期间不退化。也可尝试延长第一阶段训练至主干真正收敛后再接入通信。

再次，本文后续应进一步研究状态条件化的通信因果性约束。现有 Huber 优势引导杠杆已证明通信残差可以被训练到足以改变攻击动作 logits，而稳定通信介入方案又证明动作层暴露可以被进一步调稳；但这些改变并未稳定转化为频繁的“同伴优于本地”决策翻转。因此，后续更应关注通信何时具有因果价值，例如本地策略高不确定、队友高置信目标与本地目标冲突、或本地策略即将做出错误攻击选择的状态。相比全局提高通信影响力，这类因果状态筛选更可能把“通信路由正确”转化为“通信在该介入时介入”，并把当前“动作层暴露已修复、因果决策改变仍未解决”的边界继续向前推进。

此外，当前工作主要停留在仅攻击流的通信范围。移动通信流在本文中由于更容易出现高熵退化而未被深入修复，后续可借鉴攻击流的修复经验将其纳入。

最后，除 SMAC 外，还可将本文的通信隔离评估方法论推广到更复杂环境中，检验通信价值的环境依赖性。
